\section{Executive Summary}\label{sec:exec}

\begin{callout}[title=One-paragraph verdict]
QuarkyLab, the NetFRAME / km-cluster home lab, is an unusually mature environment
that runs three genuinely production workloads at once: DUNE neutrino-physics GPU
research, a multi-tenant student AI/ML teaching cluster, and a self-hosted LLM
inference platform. Live read-only verification on 2026-07-23 found the fabric
healthy: a 7 of 7 quorate Proxmox cluster, GPUs correctly assigned, backups fresh
(14 snapshots dated 2026-07-23), both UPS units online at full charge, and the
tracked repositories secrets-clean. The environment is well ahead of typical
home-lab practice in operational discipline. The material risks are concentrated,
well understood, and mostly already in the operator's own backlog: a single
faulted storage drive on a pool that has no second copy, a backup topology that is
effectively one-copy / one-site, and two single points of failure (OPNsense and
Randy) that the operator is already planning to address.
\end{callout}

\subsection{Overall assessment}

This is a strong environment. The header-line facts were independently confirmed
against live hardware rather than taken from documentation (E-01 through E-08).
The cluster is quorate with all corosync links connected (E-01); QuarkyLab carries
one RTX 8000 (48 GB) and Jarvis two RTX 6000 (24 GB each), which confirms the
repository and \emph{contradicts} the review specification's assertion that the
assignment was reversed (E-02, E-03, DR-01); the ZFS \code{datastore} pool is
clean and last night's backups completed (E-06, E-08). Nothing found during this
review indicates an imminent outage. The weaknesses are real but bounded, and the
operator has already produced a self-audit (commit \code{65f3681}, 2026-07-22)
that names most of them.

\subsection{Major strengths}

\begin{itemize}
  \item \textbf{Deterministic AI-operations safety.} The netframe-monitor / Jarvis
  stack is genuinely enterprise-grade in its safety design: a single deterministic
  evidence engine that every report path uses, model self-confidence stripped and
  forbidden, a three-action human-gated remediation allowlist with no autonomous
  execution, universal policy screening including the chat and Discord surfaces,
  untrusted-telemetry handling, and full survivability if the model is offline
  (E-14).
  \item \textbf{Disciplined runbook and disaster-recovery culture.} Formal RCAs and
  AARs, ADRs, an offline break-glass credential path independent of the cluster,
  a node-unreachable triage runbook, OPNsense serial-console DR with age-encrypted
  config backups, and a laptop-recovery runbook (E-21).
  \item \textbf{Real, applied network segmentation.} Seven VLANs, BMCs isolated to
  VLAN 20 with an \emph{applied} egress clamp, a VLAN 30 to VLAN 1 block, and DNS
  pinholes for isolated VLANs (E-21). This is deployed, not aspirational.
  \item \textbf{Config-as-code monitoring.} Grafana alerting, exporters, and
  textfile collectors are versioned; there are dead-man heartbeats for the pipeline
  itself (E-15).
  \item \textbf{Honest self-audit already in flight.} The operator is mid-way
  through a reconciliation branch that catalogs the documentation drift found here.
\end{itemize}

\subsection{Highest risks}

Ordering below reflects the corrections established late in the review. The
storage items were downgraded on live evidence, and the alerting item was
root-caused and closed.

\begin{enumerate}
  \item \sevhigh{} \textbf{Backups are effectively one-copy / one-site (QL-R03).}
  Every backup lives only on Randy: one PBS, one ZFS pool, one rack, one UPS bus (E-08).
  This is now the single largest unmitigated exposure.
  \item \sevhigh{} \textbf{OPNsense is a single point of failure (QL-R04).} One VM
  on pve2 provides routing, firewall, DHCP, and inter-VLAN gateway for the whole LAN.
  \item \sevhigh{} \textbf{Randy is a single storage node (QL-R05).} PBS, research
  NFS, and RKE2 persistent volumes all depend on it.
  \item \sevmed{} \textbf{A live GitHub token sits in plaintext at rest (QL-R07).}
  \code{\textasciitilde{}/github\_helper.py} contains a classic PAT; it is
  gitignored and never committed, but usable if the workstation is compromised
  (E-17). Two other plaintext secret files sit alongside it.
  \item \sevmed{} \textbf{Reduced parity on \code{bulk} (QL-R01), impact corrected.}
  \code{mpathv} is FAULTED in \code{raidz2-2}, leaving one parity drive. Originally
  rated urgent on the assumption that \code{bulk} held DUNE research data; live
  inspection (E-29) shows the pool is essentially empty, so this is loss of redundancy
  rather than imminent data loss. Replacement is pending operator hardware.
  \item \sevmed{} \textbf{\code{bulk} still has no second copy (QL-R02).} Low impact
  today because the pool is near-empty, but the exposure becomes real the moment
  research data migrates onto it, so the second copy must exist first. The Ares restic
  backup also currently targets this degraded pool.
  \item \sevmed{} \textbf{Documentation drift (QL-R19).} Recent changes landed in
  CLAUDE.md but not in README, topology, or the docs mirror; two review-spec
  "corrections" are inverted versus live reality and must not be applied. This review
  added a third: a phantom EX2300 switch that has never existed (DR-15).
\end{enumerate}

\subsection{Changes applied during this review}

The review was read-only by default. Two corrective changes were made after explicit
operator approval, and both were verified afterwards by the reviewer.

\begin{itemize}
  \item \textbf{ZFS device-level alerting deployed (QL-R11, closed).} The
  \code{ZfsPoolDegraded} rule could never have fired for \code{bulk}: it watches a
  pool-level health metric that reads ONLINE while raidz2 retains redundancy
  (\code{max\_over\_time} over 7 days = 0). No per-device metric existed. A new
  \code{node\_zfs\_pool\_nonhealthy\_devices} metric and a critical
  \code{ZfsDeviceFaulted} rule were added and deployed to Randy, QuarkyLab and Jarvis;
  verified live in Prometheus with \code{bulk} = 1 and all other pools 0.
  \item \textbf{EX3400 pinned as spanning-tree root (QL-R23, closed).} No
  \code{bridge-priority} was configured, so a UniFi access switch had won the root
  election by MAC address while documentation claimed the core was root.
  \code{bridge-priority 4096} was committed; the core is now genuinely root.
  \item \textbf{Not done: the drive replacement (QL-REC01)} remains pending operator
  hardware.
\end{itemize}

\subsection{Dimension assessments}

\begin{itemize}
  \item \textbf{Reliability:} Strong at the compute/cluster layer (7/7 quorum,
  pinned GPU stacks intact); weakened by concentrated SPOFs (OPNsense, Randy) and
  HA fencing left in standby.
  \item \textbf{Security:} Well above home-lab norm - applied segmentation, least-
  privilege API tokens, working SIEM. Residual: at-rest workstation secrets,
  unauthenticated LAN AI/registry endpoints, guest SIEM blind spot.
  \item \textbf{Recoverability:} Good documentation and an offline break-glass path;
  limited by the single backup site and one unverified restore-readiness gap
  (pve1 unbacked).
  \item \textbf{Power resilience:} Both UPS online and healthy, but UPS-A is loaded
  to 59 percent with roughly 9.5 minutes of runtime (E-07), an imbalance worth
  rebalancing.
  \item \textbf{AI-system reliability:} Excellent safety posture; the remaining lift
  is calibration/eval, not trustworthiness.
  \item \textbf{Scientific and student workloads:} Genuinely well engineered - SLURM
  shard plus per-job MPS VRAM caps plus apptainer isolation, research preemption,
  and reproducible container images.
  \item \textbf{Enterprise alignment:} Comparable-to-mature in monitoring, docs,
  IaC, and secrets discipline; material gaps in HA/redundancy and offsite DR.
\end{itemize}

\subsection{Top 10 prioritized recommendations}

\begin{enumerate}
  \item Replace the faulted \code{bulk} drive and scrub (QL-REC01, urgent).
  \item Verify the \code{ZfsPoolDegraded} alert actually reached Discord (QL-REC02, read-only).
  \item Add a second backup copy / restore the 3-2-1 principle (QL-REC03).
  \item Back up the standalone pve1 guests (QL-REC04).
  \item Rotate the at-rest GitHub PAT and the two plaintext secret files (QL-REC05).
  \item Deploy prometheus-pve-exporter for cluster/quorum/guest metrics (QL-REC06).
  \item Add the missing alert rules: PBS-stale, verify-failed, broad cert-expiry (QL-REC07).
  \item Rebalance UPS loads and confirm NUT graceful-shutdown sequencing (QL-REC08).
  \item Reconcile EX3400 SSH access and capture switch config to git via Oxidized (QL-REC09).
  \item Land the documentation reconciliation PR (D-01 through D-10) (QL-REC11).
\end{enumerate}

Building the OPNsense CARP HA pair (QL-REC10) is the highest-value medium-term
structural improvement and removes the top availability SPOF.

\subsection{Explicitly accepted risks and where no change is recommended}

The following are sound decisions that should stand: the pinned
\code{6.14.11-9-pve} kernels on the GPU nodes must not be upgraded (QL-REC20);
the raidz1 vdevs on QuarkyLab and Jarvis are an accepted single-parity tradeoff
given hot spares, clean scrubs, and low utilization (QL-R21); single-node
observability is a reasonable posture at this scale given the dead-man heartbeats
(QL-R10); and the VLAN 1 egress lockdown remains deliberately unenforced pending
its prerequisites. Stable systems should not be changed without a compelling
reason, and most of this environment is stable.

\subsection{Confidence and limitations}

Confidence in the live-verified conclusions (cluster health, GPU assignment,
storage state, backup freshness, power state, and now the switch) is \textbf{high}:
each was confirmed by direct read-only query and cross-checked against the prior
audit. The EX3400 was inspected live later in the review once the operator supplied
credentials (E-22, E-25), which lifted the original switch limitation and produced
the spanning-tree finding below. Confidence in the firewall control plane remains
\textbf{medium}: OPNsense was not queried (no credentials used, by design, E-23), so
firewall specifics still rest on repository documentation. The review was read-only
by default; the two changes listed below were made only after explicit operator
approval, and were verified afterwards. See Section~\ref{sec:limits} for the full
limitations register.
